\documentclass{article}
\usepackage[utf8]{inputenc} % Para la codificación UTF-8
\usepackage[T1]{fontenc}    % Para la codificación de la fuente
\usepackage{amsmath}        % Útil para formateo de texto
\usepackage{listings}       % Para incluir bloques de código C++
\usepackage{xcolor}         % Para usar colores en los bloques de código
\usepackage{underscore}     % Para que el carácter '_' se imprima literalmente en modo texto
\usepackage[spanish]{babel} % Para soporte de idioma español (acentos, ñ)
\selectlanguage{spanish}    % Seleccionar el idioma español
\usepackage[margin=1in]{geometry} % Ajustar márgenes para evitar Overfull \hbox y mejorar la estética

% Configuración del paquete listings para bloques de código C++
\lstset{
    language=C++,                       % Lenguaje de programación
    basicstyle=\ttfamily\small,         % Estilo de fuente (tipo máquina de escribir, pequeña)
    numbers=none,                       % No mostrar números de línea
    breaklines=true,                    % Permitir saltos de línea automáticos
    breakatwhitespace=false,            % Permitir saltos dentro de palabras largas
    frame=single,                       % Dibujar un marco alrededor del bloque de código
    captionpos=b,                       % Posición del título (bottom)
    keepspaces=true,                    % Mantener espacios
    showspaces=false,                   % No mostrar espacios como símbolos
    showstringspaces=false,             % No mostrar espacios en cadenas
    % Manejo de caracteres especiales y acentos dentro de comentarios y cadenas EN LISTINGS
    % Usamos comandos LaTeX para los caracteres dentro de las cadenas literales para máxima compatibilidad
    literate={á}{{\'a}}1 {é}{{\'e}}1 {í}{{\'i}}1 {ó}{{\'o}}1 {ú}{{\'u}}1
             {Á}{{\'A}}1 {É}{{\'E}}1 {Í}{{\'I}}1 {Ó}{{\'O}}1 {Ú}{{\'U}}1
             {ñ}{{\~n}}1 {Ñ}{{\~N}}1 {¿}{{\textquestiondown}}1 {¡}{{\textexclamdown}}1
             {ü}{{\"u}}1 {Ü}{{\"U}}1
}

% Define un comando personalizado para código/palabras clave en línea
\newcommand{\inlinecode}[1]{\texttt{#1}}

\begin{document}

\section*{Documentación de Excepciones del Sistema de Gestión de Bóvedas}

Esta sección documenta la jerarquía de excepciones personalizadas diseñada para el sistema de transporte de valores, basada en el flujo de trabajo y las clases fundamentales identificadas. Estas excepciones son cruciales para garantizar la robustez, claridad y facilidad de depuración del sistema en un entorno de misión crítica.

Todas las excepciones personalizadas heredan de una clase base común, \inlinecode{BovedaException}, que a su vez deriva de \inlinecode{std::runtime_error}, integrándose así con el sistema de manejo de excepciones estándar de \inlinecode{C++}.

\subsection*{Jerarquía de Excepciones}

La siguiente estructura muestra la relación de herencia de las clases de excepción (utilizando caracteres ASCII para compatibilidad con \inlinecode{verbatim}):
\begin{verbatim}
std::exception
  L-- std::runtime_error
        L-- BovedaException
              |-- SaldoInsuficienteException
              |-- ActivoNoDisponibleException
              |-- OperacionInvalidaException
              |-- TipoOperacionNoSoportadoException
              |-- DatosInvalidosException
              |-- BovedaNoEncontradaException
              |-- EntidadBancariaNoEncontradaException
              |-- TransportadoraNoDisponibleException
              |-- ConfiguracionInvalidaException
              L-- ErrorInternoSistemaException
\end{verbatim}

\subsection*{Clases de Excepción}

\subsubsection*{class \inlinecode{BovedaException : public std::runtime_error}}
\vspace{2mm} % Pequeño espacio vertical para separar el título del listing
\begin{lstlisting}
class BovedaException : public std::runtime_error {
public:
    explicit BovedaException(const std::string& message);
};
\end{lstlisting}
\paragraph*{Propósito:} Esta es la clase base para \textbf{todas} las excepciones personalizadas espec\'ificas del dominio del sistema de gesti\'on de b\'ovedas. Proporciona una interfaz com\'un y permite la captura gen\'erica de cualquier error relacionado con la l\'ogica de negocio de nuestro sistema.

\paragraph*{Hereda de:} \inlinecode{std::runtime_error}, lo que la hace compatible con el manejo de excepciones est\'andar de \inlinecode{C++} para errores en tiempo de ejecuci\'on.

\paragraph*{Constructor:}
\begin{itemize}
    \item \inlinecode{explicit BovedaException(const std::string\& message)}
    \begin{itemize}
        \item \textbf{Parámetro \inlinecode{message}:} Un mensaje descriptivo del error que ocurri\'o. Este mensaje se puede recuperar con el m\'etodo \inlinecode{what()}.
    \end{itemize}
\end{itemize}

\paragraph*{Uso Típico:} No se lanza directamente. Su funci\'on principal es servir como base para otras excepciones m\'as espec\'ificas y permitir la captura polim\'orfica de errores de la b\'oveda.

\subsubsection*{Excepciones Relacionadas con el Estado y los Fondos}

\paragraph*{\textbf{class \inlinecode{SaldoInsuficienteException : public BovedaException}}}
\vspace{2mm} % Pequeño espacio vertical
\begin{lstlisting}
class SaldoInsuficienteException : public BovedaException {
public:
    explicit SaldoInsuficienteException(const std::string& message = "Error de fondos: Saldo insuficiente en la boveda de origen.");
};
\end{lstlisting}
\paragraph*{Propósito:} Indica que una operaci\'on no puede completarse debido a que la b\'oveda de origen (o la entidad bancaria) no posee la cantidad de fondos requerida.

\paragraph*{Escenario Típico:}
\begin{itemize}
    \item Un intento de enviar S/ 1 mill\'on con solo S/ 800,000 disponibles.
\end{itemize}

\paragraph*{\textbf{class \inlinecode{ActivoNoDisponibleException : public BovedaException}}}
\vspace{2mm} % Pequeño espacio vertical
\begin{lstlisting}
class ActivoNoDisponibleException : public BovedaException {
public:
    explicit ActivoNoDisponibleException(const std::string& message = "Error de activo: El tipo de activo solicitado no esta disponible o es insuficiente.");
};
\end{lstlisting}
\paragraph*{Propósito:} Se lanza si se intenta realizar una operaci\'on con un tipo de activo (moneda, bonos, joyas) que la b\'oveda de origen no posee, o no en la cantidad necesaria, incluso si el saldo total de la b\'oveda fuera suficiente con otros activos.

\paragraph*{Escenario Típico:}
\begin{itemize}
    \item Intentar retirar 1,000 USD de una b\'oveda que solo almacena Soles Peruanos (PEN).
\end{itemize}

\subsubsection*{Excepciones Relacionadas con la Validación de Datos y la Lógica de Negocio}

\paragraph*{\textbf{class \inlinecode{OperacionInvalidaException : public BovedaException}}}
\vspace{2mm} % Pequeño espacio vertical
\begin{lstlisting}
class OperacionInvalidaException : public BovedaException {
public:
    explicit OperacionInvalidaException(const std::string& message = "Error de operacion: La operacion es invalida debido a reglas de negocio o estado.");
};
\end{lstlisting}
\paragraph*{Propósito:} Indica una violaci\'on de las reglas de negocio o una transici\'on de estado ilegal para una operaci\'on. Es una excepci\'on general para errores de l\'ogica de negocio que no encajan en otras categor\'ias m\'as espec\'ificas.

\paragraph*{Escenarios Típicos:}
\begin{itemize}
    \item Intentar cancelar una operaci\'on que ya ha sido marcada como "Completada" o "Cancelada".
    \item Intentar crear una operaci\'on \inlinecode{PROVISION\_BCRP} (recibir fondos del BCRP) desde una b\'oveda comercial, cuando esta operaci\'on solo es v\'alida para el BCRP.
    \item Intentar una transferencia interna con un monto que excede un l\'imite predefinido sin la aprobaci\'on adecuada.
\end{itemize}

\paragraph*{\textbf{class \inlinecode{TipoOperacionNoSoportadoException : public BovedaException}}}
\vspace{2mm} % Pequeño espacio vertical
\begin{lstlisting}
class TipoOperacionNoSoportadoException : public BovedaException {
public:
    explicit TipoOperacionNoSoportadoException(const std::string& message = "Error de tipo: El tipo de operacion especificado no es soportado.");
};
\end{lstlisting}
\paragraph*{Propósito:} Se lanza si el sistema recibe o intenta procesar un tipo de operaci\'on que no est\'a definido o no es reconocido por el sistema (ej. un valor de enumeraci\'on inv\'alido para \inlinecode{TipoOperacion}).

\paragraph*{Escenario Típico:}
\begin{itemize}
    \item Una integraci\'on externa intenta crear una transacci\'on con un \inlinecode{TipoOperacion} como "AJUSTE\_INVENTARIO" que no existe en el cat\'alogo de operaciones del sistema.
\end{itemize}

\paragraph*{\textbf{class \inlinecode{DatosInvalidosException : public BovedaException}}}
\vspace{2mm} % Pequeño espacio vertical
\begin{lstlisting}
class DatosInvalidosException : public BovedaException {
public:
    explicit DatosInvalidosException(const std::string& message = "Error de datos: Los datos de entrada son invalidos o tienen formato incorrecto.");
};
\end{lstlisting}
\paragraph*{Propósito:} Indica que los datos proporcionados para una operaci\'on o entidad son inv\'alidos, tienen un formato incorrecto o no cumplen con restricciones b\'asicas de entrada.

\paragraph*{Escenario Típico:}
\begin{itemize}
    \item Intentar crear una operaci\'on con un monto negativo o cero.
    \item Proporcionar un ID de b\'oveda no num\'erico o una fecha en un formato incorrecto.
\end{itemize}

\subsubsection*{Excepciones Relacionadas con la Disponibilidad y Existencia de Entidades}

\paragraph*{\textbf{class \inlinecode{BovedaNoEncontradaException : public BovedaException}}}
\vspace{2mm} % Pequeño espacio vertical
\begin{lstlisting}
class BovedaNoEncontradaException : public BovedaException {
public:
    explicit BovedaNoEncontradaException(const std::string& message = "Error de entidad: La boveda especificada no fue encontrada.");
};
\end{lstlisting}
\paragraph*{Propósito:} Se lanza cuando el sistema intenta interactuar con una b\'oveda (ya sea de origen o destino) cuyo identificador no corresponde a ninguna b\'oveda registrada en el sistema.

\paragraph*{Escenario Típico:}
\begin{itemize}
    \item Ingresar un ID de b\'oveda (\inlinecode{BVD\_XYZ}) incorrecto o inexistente al programar una transferencia.
\end{itemize}

\paragraph*{\textbf{class \inlinecode{EntidadBancariaNoEncontradaException : public BovedaException}}}
\vspace{2mm} % Pequeño espacio vertical
\begin{lstlisting}
class EntidadBancariaNoEncontradaException : public BovedaException {
public:
    explicit EntidadBancariaNoEncontradaException(const std::string& message = "Error de entidad: La entidad bancaria especificada no fue encontrada.");
};
\end{lstlisting}
\paragraph*{Propósito:} Indica que una operaci\'on hace referencia a una entidad bancaria (un banco o el BCRP) que no est\'a registrada o no existe en el sistema.

\paragraph*{Escenario Típico:}
\begin{itemize}
    \item Intentar una \inlinecode{TRASLADO\_INTERBANCARIO} a un banco con un RUC que no corresponde a ninguna entidad bancaria conocida.
\end{itemize}

\paragraph*{\textbf{class \inlinecode{TransportadoraNoDisponibleException : public BovedaException}}}
\vspace{2mm} % Pequeño espacio vertical
\begin{lstlisting}
class TransportadoraNoDisponibleException : public BovedaException {
public:
    explicit TransportadoraNoDisponibleException(const std::string& message = "Error de servicio: La transportadora no esta disponible o no cubre la zona.");
};
\end{lstlisting}
\paragraph*{Propósito:} Se lanza si la empresa transportadora de valores asignada a una operaci\'on no est\'a disponible para el servicio, no cubre la ubicaci\'on de destino, o no est\'a habilitada en el sistema.

\paragraph*{Escenario Típico:}
\begin{itemize}
    \item Intentar programar un env\'io con "Transportes XYZ" a una ubicaci\'on remota que no est\'a dentro de su \'area de servicio.
    \item La transportadora seleccionada est\'a marcada como "inactiva" en el sistema.
\end{itemize}

\subsubsection*{Excepciones de Sistema / Internas}

\paragraph*{\textbf{class \inlinecode{ConfiguracionInvalidaException : public BovedaException}}}
\vspace{2mm} % Pequeño espacio vertical
\begin{lstlisting}
class ConfiguracionInvalidaException : public BovedaException {
public:
    explicit ConfiguracionInvalidaException(const std::string& message = "Error de configuracion: El sistema no esta configurado correctamente.");
};
\end{lstlisting}
\paragraph*{Propósito:} Indica un problema con la configuraci\'on interna del sistema que impide el correcto funcionamiento. Estos errores suelen ser detectados al inicio o durante la inicializaci\'on de m\'odulos.

\paragraph*{Escenario Típico:}
\begin{itemize}
    \item Los identificadores predefinidos para el BCRP o ciertas b\'ovedas clave no han sido cargados o son incorrectos en la configuraci\'on.
\end{itemize}

\paragraph*{\textbf{class \inlinecode{ErrorInternoSistemaException : public BovedaException}}}
\vspace{2mm} % Pequeño espacio vertical
\begin{lstlisting}
class ErrorInternoSistemaException : public BovedaException {
public:
    explicit ErrorInternoSistemaException(const std::string& message = "Error interno: Se ha producido un fallo inesperado en el sistema.");
};
\end{lstlisting}
\paragraph*{Propósito:} Una excepci\'on de "\'ultimo recurso" para errores inesperados y no clasificados que indican un fallo en la l\'ogica interna del sistema o una condici\'on que no deber\'ia ocurrir bajo operaci\'on normal.

\paragraph*{Escenario Típico:}
\begin{itemize}
    \item Un fallo en una dependencia interna, una corrupci\'on de datos en memoria o un estado de objeto inconsistente que no puede ser atribuido a una excepci\'on m\'as espec\'ifica.
\end{itemize}

\subsection*{Principios de Uso de las Excepciones}

\paragraph*{\textbf{Lanzar (\inlinecode{throw})}}
Cuando se detecta una condici\'on de error que la funci\'on actual no puede resolver, se crea una instancia de la excepci\'on m\'as espec\'ifica y se lanza utilizando la palabra clave \inlinecode{throw}.
\begin{lstlisting}
throw SaldoInsuficienteException("La cuenta de origen no tiene fondos suficientes.");
\end{lstlisting}

\paragraph*{\textbf{Capturar (\inlinecode{catch})}}
Se utilizan bloques \inlinecode{try-catch} para interceptar las excepciones lanzadas. Es crucial colocar los bloques \inlinecode{catch} de las excepciones m\'as espec\'ificas antes que los m\'as gen\'ericos en una secuencia \inlinecode{try-catch}.
\begin{lstlisting}
try {
    // Código que puede lanzar una excepción
    realizarOperacion();
} catch (const SaldoInsuficienteException& e) {
    // Tratamiento específico para saldo insuficiente
    std::cerr << "ERROR DE SALDO: " << e.what() << std::endl;
} catch (const BovedaException& e) {
    // Tratamiento genérico para cualquier error del sistema de bóveda
    std::cerr << "ERROR DEL SISTEMA BOVEDA: " << e.what() << std::endl;
} catch (const std::exception& e) {
    // Tratamiento para cualquier otra excepción estándar
    std::cerr << "ERROR INESPERADO: " << e.what() << std::endl;
}
\end{lstlisting}

\paragraph*{\textbf{Mensajes Descriptivos (\inlinecode{what()})}}
Todas las clases de excepci\'on heredan el m\'etodo \inlinecode{what()} de \inlinecode{std::runtime_error} (o \inlinecode{std::exception}). Este m\'etodo devuelve la cadena de caracteres (el mensaje) que se pas\'o al constructor de la excepci\'on, proporcionando informaci\'on vital sobre la causa del error.

\paragraph*{\textbf{Relanzamiento (\inlinecode{throw;})}}
En algunos casos, una funci\'on puede capturar una excepci\'on, realizar un procesamiento parcial (ej. loguear el error, limpiar recursos locales) y luego relanzarla (\inlinecode{throw;}) para que un nivel superior de la aplicaci\'on pueda manejarla de forma m\'as global.
\begin{lstlisting}
try {
    // Código que lanza excepción
    throw SaldoInsuficienteException("Faltan fondos.");
} catch (const BovedaException& e) {
    // Primer tratamiento: loguear el error
    std::cerr << "LOG: Error capturado internamente: " << e.what() << std::endl;
    // Relanzar para tratamiento superior
    throw;
}
\end{lstlisting}

\end{document}